\section{Every irrational angle has cubic scales}\label{sec:irrational}
We keep the two-letter alphabet $\{I,R_{2\pi\alpha}\}$ and the planar resource and error conventions. No Diophantine lower bound is assumed in this section. The needed continued-fraction estimates are classical \cite{Cassels}; we include their short proof to specify the exact finite inequalities being used.

\begin{lemma}[Convergent separation]\label{lem:cf}
Let $p_n/q_n$ be the convergents of an irrational $\alpha$, ignoring the finitely many initial indices with repeated denominator. Then
\begin{equation}\label{eq:cf-estimates}
 \frac1{q_{n+1}+q_n}<|q_n\alpha-p_n|<\frac1{q_{n+1}},
 \qquad
 \min_{1\le l<q_n}\norm{l\alpha}_{\R/\Z}>
                  \frac1{q_n+q_{n-1}}>\frac1{2q_n}.
\end{equation}
\end{lemma}
\begin{proof}
The convergent recurrences imply
$p_nq_{n-1}-p_{n-1}q_n=\pm1$ and alternating signs of $e_n=q_n\alpha-p_n$. If $\xi$ is the next complete quotient, then
\[
 \alpha=\frac{p_n\xi+p_{n-1}}{q_n\xi+q_{n-1}},\qquad
 a_{n+1}<\xi<a_{n+1}+1.
\]
It follows that $|e_n|=1/(q_n\xi+q_{n-1})$, proving the first pair.

For $0<l<q_n$ and an integer $p$, the determinant identity writes
$(p,l)=a(p_n,q_n)+b(p_{n-1},q_{n-1})$ with integers $a,b$. Here $b\ne0$; $a,b$ cannot both be positive, and cannot both be negative. If $a=0$ the error has modulus at least $|e_{n-1}|$. Otherwise they have opposite signs, so $ae_n$ and $be_{n-1}$ have the same sign and again
$|l\alpha-p|\ge|e_{n-1}|$. Minimize over $p$ and apply the first pair with index $n-1$. This proves the remaining inequalities.
\end{proof}

\begin{theorem}[Universal cubic subsequences]\label{thm:cubic-scales}
Set $\kappa=\rho/\sqrt2-2\varepsilon>0$, $C_\kappa=8(1+\log(1/\kappa))$, and
\[
 M_n=\lceil C_\kappa q_n^3\rceil,\qquad
 m_\kappa=\lceil100(C_\kappa+1)\rceil.
\]
For all sufficiently large $n$,
\begin{equation}\label{eq:cubic-scales}
               q_n/2<W_{M_n,\varepsilon}(\alpha)\le m_\kappa q_n.
\end{equation}
The upper bound is exact, uses common command rows and requires no regularity of the partial quotients.
\end{theorem}
\begin{proof}
Use words of length $B=q_n-1$ whose number of initial active letters ranges uniformly over $0,\ldots,q_n-1$. Lemma~\ref{lem:cf} gives $q_n$ distinct phases with separation greater than $1/(2q_n)$. If the proposed peak $K$ were at most $q_n/2$, Corollary~\ref{cor:separated} and the packet budget would give
\[
 \left\lfloor\frac{M_n}{q_n-1}\right\rfloor\frac1{4q_n^2}
                           \le\log(1/\kappa).
\]
The left side is at least $C_\kappa/4-1/(4q_n^2)$, strictly greater than the right side. This proves the lower bound against all hidden-state machines.

For the upper bound use the polygon denominator $Q_n=m_\kappa q_n$. Distance to the nearest integer is subadditive, so
\[
 \norm{Q_n\alpha}_{\R/\Z}
 \le m_\kappa|q_n\alpha-p_n|<m_\kappa/q_{n+1}\le m_\kappa/q_n.
\]
Consequently
$Q_n^2/[100\norm{Q_n\alpha}_{\R/\Z}]>m_\kappa q_n^3/100
\ge(C_\kappa+1)q_n^3\ge M_n$.
Equation~\eqref{eq:eta-upper} gives the exact upper construction. The constants have been chosen for transparency, not optimality.
\end{proof}

\begin{proposition}[Resonant long horizons]\label{prop:resonant-horizons}
For any irrational $\alpha$ and integer $q\ge4$ with
$N_q=\lfloor q^2/(100\norm{q\alpha}_{\R/\Z})\rfloor\ge1$, one has $W_{N_q,0}(\alpha)\le q$. If $\alpha$ is Liouville, there is a sequence $N\to\infty$ on which
$\log W_{N,\varepsilon}(\alpha)/\log N\to0$.
\end{proposition}
\begin{proof}
The first assertion is the finite enclosure implication \eqref{eq:eta-upper}. For a Liouville number choose arbitrarily large $A$ and unbounded denominators $q$ with
$0<\norm{q\alpha}_{\R/\Z}<q^{-A}$. The first assertion then supplies horizons $N_q\ge q^{A+2}/200$ for large $q$, with width at most $q$. The logarithmic ratio is at most $1/(A+2)+o(1)$. A diagonal sequence with $A\to\infty$ proves the claim, since widths are at least two.
\end{proof}

\begin{proof}[Proof of Theorem~\ref{thm:liouville-main}]
Equation~\eqref{eq:cubic-scales} and $M_n=\Theta(q_n^3)$ give the first assertion and a subsequence of logarithmic exponent $1/3$. Proposition~\ref{prop:resonant-horizons} gives the zero lower limit for Liouville angles, and \eqref{eq:sqrt-upper} gives the upper limit at most $1/2$.

It remains to verify unboundedness, which does not follow from a subsequence alone. Fix any proposed width $K$. The $2K$ phases of counts $0,\ldots,2K-1$ are distinct for an irrational angle and have some positive minimum separation $s_K$. Independent words of length $2K-1$ imply
$\lfloor N/(2K-1)\rfloor s_K^2\le\log(1/\kappa)$ for a width-$K$ machine. This fails for all sufficiently large $N$. Thus $W_{N,\varepsilon}\to\infty$.
\end{proof}

\begin{corollary}[An explicit fluctuating alphabet]\label{cor:liouville-explicit}
Let $\alpha_L=\sum_{j=1}^{\infty}10^{-j!}$. For $q_m=10^{m!}$ and all large $m$, the horizon
\[
                     N_m=\lfloor q_m^{m+2}/200\rfloor
\]
admits an exact machine with at most $q_m$ labels. The same fixed two-letter alphabet also has the cubic subsequences of Theorem~\ref{thm:cubic-scales}; hence no single logarithmic width exponent exists. Adding reflection preserves these conclusions.
\end{corollary}
\begin{proof}
The first $m$ summands have denominator dividing $q_m$, and their remaining positive tail is less than $2\cdot10^{-(m+1)!}$. Thus
$0<\norm{q_m\alpha_L}_{\R/\Z}<2q_m^{-m}$ for large $m$, proving the horizon bound by \eqref{eq:eta-upper}. The same tail estimates preclude rationality and show the Liouville property. The general cubic theorem applies. Reflection leaves the upper polygons invariant, while the rotation-only test words give the same lower bounds.
\end{proof}

These theorems describe different scales of one and the same experiment. A small label count on a long resonant horizon is not a constant-width simulator for all horizons, and it need not persist across the gaps between successive approximants. The exact finite optimum and the Liouville upper logarithmic limit remain separate quantitative questions. None of the proofs substitutes approximate output for the exact upper realizations.
